\section{Similarity-benchmark protocol}
\label{app:benchmarks}

The field-level benchmarks cover Economics, Psychology, and Physics, with data described in Appendix~\ref{app:data}.
For \doctolora, we evaluate embeddings produced by the Qwen3-4B checkpoint across the primary benchmark, including all three field tasks and the five S2AND datasets.
For retrieval, we mean-pool the adapter tensors across the rank dimension as defined in \secref{doc2lora-embedding}.
Because the transformation $g_\theta$ is linear, pooling before or after applying $g_\theta$ yields identical vectors.
Text baselines comprise \texttt{SPECTER2}~\citep{singh2023scirepeval}, \texttt{SBERT} (all-mpnet-base-v2)~\citep{reimers2019sbert}, \texttt{Instructor} (instructor-large)~\citep{su2023instructor}, \texttt{EmbeddingGemma} (\texttt{embeddinggemma-300m})~\citep{embeddinggemma2025}, and \texttt{GTE} (\texttt{gte-large-en-v1.5})~\citep{li2023gte} used by \texttt{T2L}.
All text encoders run with default settings without task-specific tuning.
\tabref{similarity} lists the absolute scores, and the mean ranks cited in \secref{similarity} are computed from that table.

For the S2AND datasets, we evaluate against the earlier \texttt{SPECTER} model because the S2AND system is designed around those representations.
The benchmark does not provide \texttt{SPECTER2} vectors.
Generating vectors locally would require changing both the text preprocessing and the encoder, which causes the evaluation to diverge from the published benchmark.
Because official vectors do not exist for the other text baselines, we compute their embeddings ourselves from the published titles and abstracts.

For next-paper prediction, papers are split per author over time.
Each query paper is paired with the next paper written by the same author, which forms a positive pair. 
We also pair the query paper with a randomly sampled paper that the author did not write, which forms a negative pair.
For topic classification, we assign OpenAlex papers to one of 26 Scopus or ASJC classes, APS papers to one of 18 Physical Review codes, and sample 60000 papers seen by all encoders, splitting into 80 percent training and 20 percent test, not based on year.
Collaboration prediction uses anchor years 2000, 2004, 2008 for APS and 2008, 2012, 2016 for Economics and Psychology
At each anchor, an author is represented by averaging their papers from the previous four years
Candidate pairs are author pairs at network distance two who share at least one coauthor and have no joint publication.
A positive example is a candidate pair who coauthor a paper together in the next three years after the anchor, while a negative is a pair who do not coauthor in that period.
Because positive and negative pairs have identical local coauthorship structure, only content similarity can distinguish them.
All candidate pairs from all anchor windows are pooled to compute mean AUC-ROC.
For author-name disambiguation, we use the five S2AND subsets from Appendix~\ref{app:data}; each signature is represented by the embedding of its paper
Signatures are clustered by agglomerative average linkage with a cosine-distance threshold tuned on training data and performance is reported using B3 F1

Uncertainties come from one thousand bootstrap resamples of per-unit evaluation scores, and the $\pm$ in \tabref{similarity} is the standard deviation of a benchmark's own resample distribution. Each method is resampled independently, so a marginal interval reflects only that method and comparing two of them is conservative. The mean ranks quoted in \secref{similarity} average over all fourteen benchmarks rather than over the scores themselves, because the benchmarks carry different metrics and sit at different levels. Resampling units depend on the task. Next-paper prediction resamples queries in clusters. Collaboration prediction resamples candidate author pairs. Topic classification resamples test papers. Disambiguation resamples signatures within name blocks by drawing B\textsuperscript{3} precision and recall jointly and then computing F1.

Relative model rankings in \tabref{similarity} remain stable across all three base models (\tabref{encoder-matrix} in Appendix~\ref{app:tables}).
The smallest model, Gemma-2-2b, trails Qwen3-4B by up to ${\sim}.05$ on the field tasks but exceeds it on Physics collaboration ($.823$ against $.795$).
Mistral-7B performs comparably to Qwen3-4B and outperforms it on Physics next-paper prediction ($.965$).
Each cell in \tabref{encoder-matrix} reports raw and +$g_\theta$ scores using the evaluation pools, metrics, and bootstrap procedures from \tabref{similarity}, omitting standard deviations for readability.
Because text baseline scores are independent of the base model, we report them only in \tabref{similarity}.
